\renewcommand{\Anno}{1996}

\renewcommand{\NumeroProblema}{4}

\renewcommand{\Testo}{
Determinare le ultime cinque cifre (cio\`{e} quelle delle unit\`{a},
decine, centinaia, migliaia e decine di migliaia) del numero
$$5^{5^5}.$$}

\renewcommand{\Risposta}{$03125$}

\renewcommand{\Aiutino}{
L'idea fondamentale \`{e} che le ultime 5 cifre 
del numero dato coincidono con le ultime 5 cifre di $5^{5}$. Per 
dimostrarlo, vi sono vari approcci possibili.
\begin{itemize}
	\item  Dimostrare che $5^{3125}-5^{5}$ \`{e} multiplo di $2^{5}$ e 
	di $5^{5}$, dunque anche di $10^{5}$. Per dimostrarlo, basta 
	fattorizzare $5^{3120}-1$ come prodotto di 4 fattori pari, di cui 
	almeno uno \`{e} multiplo di 4.

	\item Osservare che, a partire da $5^{5}$, le ultime 5 cifre delle
	potenze di 5 si ripetono con periodo 8. Questo richiede il calcolo 
	delle ultime 5 cifre di tutte le potenze di 5 fino alla tredicesima, 
	cosa che si semplifica notevolmente notando che le ultime 5 cifre di 
	$5k$ non dipendono da tutta la rappresentazione decimale del numero 
	$k$, ma solo dalle sue ultime 5 cifre.

	\item  Dimostrare che le potenze di 5 hanno periodo 8 modulo 32.
\end{itemize}
}

\renewcommand{\SoluzioneUfficiale}{
Dimostriamo che le ultime 5 cifre del numero dato coincidono con le
ultime 5 cifre di $5^{5}$, e sono quindi $03125$.  Questo \`{e}
equivalente a dimostrare che il numero $N=5^{3125}-5^{5}$ \`{e}
divisibile per $10^{5}$, il che a sua volta si riduce a dimostrare che
$N$ \`{e} divisibile sia per $5^{5}$ sia per $2^{5}$.

Scrivendo 
$$N=5^{5}\left(5^{3120}-1\right)$$
la divisibilit\`{a} per $5^{5}$ \`{e} immediata.  Ci resta quindi da
dimostrare che il secondo fattore \`{e} divisibile per $2^{5}$.
Applicando pi\`{u} volte la formula per la differenza di quadrati, si 
ha che
\begin{eqnarray*}
	5^{3120}-1 & = & (5^{1560}+1)(5^{1560}-1)\ =\ 
	(5^{1560}+1)(5^{780}+1)(5^{780}-1)\ =\ \ldots  \\
	 & = & (5^{1560}+1)(5^{780}+1)(5^{390}+1)(5^{195}+1)(5^{195}-1).
\end{eqnarray*}

Poich\'{e} i 5 fattori sono tutti pari (in quanto somma o differenza di
numeri dispari), ne segue che $5^{3120}-1$ \`{e} multiplo di $2^{5}$,
e questo completa la dimostrazione.  }

\renewcommand{\SoluzioniAlternative}{\TitoloSoluzioneAlternativaA
Determinare le ultime 5 cifre di un numero equivale a determinare la 
sua classe di congruenza modulo $10^{5}$, il che per il teorema cinese 
equivale a determinare la sua classe di congruenza modulo $5^{5}$ e 
modulo $2^{5}$.

Ora \`{e} chiaro che
\begin{equation}
	5^{3125}\equiv 5^{5}\equiv 0\pmod{5^{5}},
	\label{eq:mod1}
\end{equation}
cos\`{\i} come avviene per ogni potenza di 5 con esponente maggiore 
od uguale a 5.

Passando alla congruenza modulo $2^{5}=32$, si ha che
$$5^{8}\equiv 1\pmod{32},$$
come si pu\`{o} dimostrare con una verifica diretta o pi\`{u} 
generalmente ricordando che modulo 32 l'ordine moltiplicativo di ogni 
numero dispari \`{e} un divisore di 8. Ne segue che
\begin{equation}
	5^{3125}\equiv 5^{3120}\cdot 5^{5}\equiv (5^{8})^{390}\cdot 
	5^{5}\equiv 5^{5}\pmod{2^{5}}.
	\label{eq:mod2}
\end{equation}

Dalla (\ref{eq:mod1}) e dalla (\ref{eq:mod2}) segue quindi che
$$5^{3125}\equiv 5^{5}\equiv 3125\pmod{10^{5}},$$
da cui il risultato cercato.

\TitoloSoluzioneAlternativaB
Calcoliamo le ultime 5 cifre delle prime potenze di 5.
\begin{center}
	\renewcommand{\arraystretch}{1.3}
	\begin{tabular}{|c|c|}
		\hline
		Potenza & Ultime 5 cifre  \\
		\hline\hline
		$5^{1}$ & 00005  \\
		\hline
		$5^{2}$ & 00025  \\
		\hline
		$5^{3}$ & 00125  \\
		\hline
		$5^{4}$ & 00625  \\
		\hline
		$5^{5}$ & 03125  \\
		\hline
		$5^{6}$ & 15625  \\
		\hline
		$5^{7}$ & 78125  \\
		\hline
		$5^{8}$ & 90625  \\
		\hline
		$5^{9}$ & 53125  \\
		\hline
		$5^{10}$ & 65625  \\
		\hline
		$5^{11}$ & 28125  \\
		\hline
		$5^{12}$ & 40625  \\
		\hline
		$5^{13}$ & 03125  \\
		\hline
	\end{tabular}
\end{center}

Nel compilare la tabella non abbiamo dovuto calcolare tutte le 
potenze, ma abbiamo seplicemente osservato che le ultime 5 cifre del 
prodotto $5k$ non dipendono dall'intera scrittura decimale di $k$, ma 
solo dalle sue ultime 5 cifre. In particolare, le 5 cifre scritte 
nella seconda colonna di ogni riga si ottengono moltiplicando per 5 
le 5 cifre della riga precedente e considerando a sua volta le ultime 
5 cifre del prodotto.

Pertanto, avendo ottenuto che le ultime 5 cifre di $5^{13}$ sono le
stesse di $5^{5}$, avremo che anche le ultime 5 cifre di $5^{14}$ e
$5^{6}$ sono le stesse, e cos\`{\i} via.  Questo sostanzialmente vuol
dire che le ultime 5 cifre si ripetono con periodo $13-5=8$ a partire
da quelle di $5^{5}$ (cos\`{\i} come si pu\`{o} notare che le ultime 3
cifre si ripetono con periodo 2 e le ultime 4 cifre si ripetono con
periodo 4).

Ne segue che $5^{8k+5}$ ha le stesse 5 cifre finali di $5^{5}$ per 
ogni intero positivo, dunque in particolare $5^{3125}=5^{8\cdot 
390+5}$ termina con le stesse 5 cifre di $5^{5}$, quindi con $03125$.
}

\renewcommand{\Osservazioni}{\TitoloOsservazione La tecnica utilizzata
nella prima soluzione permette in realt\`{a} di dimostrare che tutti i
numeri della forma $5^{8k}-1$ sono multipli di $2^{5}$.  Infatti nella
scomposizione 
$$5^{8k}-1 =(5^{4k}+1)(5^{2k}+1)(5^{k}+1)(5^{k}-1)$$
abbiamo 4 fattori pari (e questo assicura la divisibilit\`{a} per 
$2^{4}$); inoltre la differenza tra gli ultimi 2 fattori \`{e} 2, e 
questo ci dice che uno di loro \`{e} divisibile per 4 e ci permette 
di guadagnare il fattore 2 mancante. }

\renewcommand{\Commenti}{\TitoloCommenti 
La prima soluzione usa solo scomposizioni elementari. Tuttavia l'idea 
fondamentale, e cio\`{e} che le ultime 5 cifre di $5^{3125}$ e $5^{5}$ 
sono le stesse, ``arriva un po' dall'alto''. Per questo motivo \`{e} 
piuttosto improbabile che un solutore non esperto arrivi di primo 
acchito ad un approccio del genere.

La seconda soluzione \`{e} invece piuttosto standard per chi ha
\emph{buona} familiarit\`{a} con le classi di resto, ma proprio per 
questo rischia di ``tagliare le gambe'' agli altri.

La terza soluzione, di sicuro meno brillante dal punto di vista 
estetico, non prerequisisce nessuna conoscenza specifica, ma arriva 
``sperimentalmente'' al punto fondamentale, cio\`{e} alla periodicit\`{a} 
delle ultime 5 cifre.
}